\subsection{庞加莱群的基数}

命题 9. --- 设 X 是一个可解圈拓扑空间，W 是 X 的拓扑的一个基。对 X 的每一点 x，群 \(\pi_{1}(X,x)\) 的基数以 sup(Card(W),Card(N)) 为上界。特别地，可解圈的可数型度量空间的庞加莱群是可数的。

事实上，空间 \(\lambda_{x}(\mathrm{X})\)——赋以终点映射——是 X 的一个连通覆叠，它在 x 处的纤维是 \(\pi_{1}(\mathrm{X}, x)\)。断言于是由 I, p. 40, 定理 3 得出。

引理 4. --- 设 X 是贝尔空间，G 是连续地作用在 X 上的一个拓扑群，B 是 X 的一个不贫的可达子集。由点 \(g \in G\)（满足 \(B \cap gB \neq \emptyset\)）所成的集合是 G 的单位元的一个邻域。

由于 B 可达，\(\complement B\) 也可达，因为 X 的可达子集的全体是一个 \(\sigma\)-代数。于是存在 X 的一个开子集 U，使 \(U \cap \complement B\) 与 \(B \cap \complement U\) 都在 X 中贫。

设 V 是 G 的单位元的一个邻域，W 是 X 的一个包含于 U 的非空开子集，使 \(V \cdot W \subset U\)。对所有 \(g \in V\)，\(U \cap gU\) 非空。

设 \(g \in G\) 使 \(B \cap gB\) 为空。关系式

\[
\begin{array}{l} \mathrm{U} \cap g \mathrm{U} = (\mathrm{U} \cap g \mathrm{U}) \cap \mathbb {C} (\mathrm{B} \cap g \mathrm{B}) \\ \qquad = (\mathrm{U} \cap g \mathrm{U}) \cap (\mathbb {C} \mathrm{B} \cup g \mathbb {C} \mathrm{B}) \\ \qquad = (\mathrm{U} \cap g \mathrm{U} \cap \mathbb {C} \mathrm{B}) \cup (\mathrm{U} \cap g \mathrm{U} \cap g \mathbb {C} \mathrm{B}) \\ \qquad \subset (\mathrm{U} \cap \mathbb {C} \mathrm{B}) \cup g (\mathrm{U} \cap \mathbb {C} \mathrm{B}) \end{array}
\]

蕴含 U ∩ gU 在 X 中贫。由于它是一个开子集且 X 是贝尔空间，它为空。于是有 \(g \notin V\)，引理由此得证。

定理 2（Shelah \(^{(2)}\)）. --- 设 X 是一个连通且局部道路连通的波兰空间，x 是 X 的一点。若 X 不可解圈，则群 \(\pi_{1}(X,x)\) 具有连续统的基数。

设 d 是定义 X 的拓扑且使 X 完备的一个度量。假设 X 不可解圈。则存在一点 \(a \in X\)，且对每个整数 \(n \geqslant 0\)，存在 X 中的一个圈 \(c_{n}\)，其像的直径 \(\leqslant 2^{-n}\)，且其类在 \(\pi_{1}(X, a)\) 中非平凡。

设 K 表示集合 \(\{0,1\}^{\mathbf{N}}\)，并给它赋以乘积拓扑——空间 \(\{0,1\}\) 赋以离散拓扑。对 K 的每个元素 \(\varepsilon = (\varepsilon_n)\)，设 \(c_{\varepsilon}\) 是由 \(c_{\varepsilon}(0) = a\) 定义的、从 I 到 X 中的映射，且对 \(2^{-n - 1} \leqslant t \leqslant 2^{-n}\)，有 \(c_{\varepsilon}(t) = c_n(2^{n + 1}t - 1)\)（若 \(\varepsilon_n = 1\)），否则 \(c_{\varepsilon}(t) = a\)。我们有 \(c_{\varepsilon}(0) = c_{\varepsilon}(1) = a\)。映射 \(c_{\varepsilon}\) 在每一点 \(t \in ]0,1]\) 处连续。它在 0 处也连续，因为有 \(d(a,c_{\varepsilon}(t)) \leqslant 2^{-n}\)（若 \(t \in [0,2^{-n}] \)）。于是映射 \(c_{\varepsilon}\) 是 X 中以 \(a\) 为基点的一个圈。

若 \(\varepsilon\) 与 \(\varepsilon'\) 是 K 的元素，满足 \((\varepsilon_0, \ldots, \varepsilon_n) = (\varepsilon_0', \ldots, \varepsilon_n')\)，则 \(c_{\varepsilon}(t) = c_{\varepsilon'}(t)\)（对所有 \(t \in [2^{-n-1}, 1]\)），且有 \(d(c_{\varepsilon}(t), c_{\varepsilon'}(t)) \leqslant d(a, c_{\varepsilon}(t)) + d(a, c_{\varepsilon'}(t)) \leqslant 2^{-n}\)（若 \(t \in [0, 2^{-n-1}]\)）。由此得出，映射 \(\varepsilon \mapsto c_{\varepsilon}\)（从 K 到空间 \(\Omega_a(X)\) 中）是连续的，当空间 \(\Omega_a(X)\) 赋以紧收敛拓扑时。

记 \(\Gamma \subset \mathrm{K} \times \mathrm{K}\) 为由满足下列条件的偶 \((\varepsilon, \varepsilon')\) 所成的集合：\(c_{\varepsilon}\) 严格同伦于 \(c_{\varepsilon'}\)。它是 K 上一个等价关系 R 的图象。

引理 5. --- 集合 \(\Gamma\) 是 K × K 的一个贫子集。

设 Z 表示空间 \(K \times K \times \mathcal{C}_{c}(\mathbf{I} \times \mathbf{I}; X)\)。空间 \(\mathcal{C}_{c}(\mathbf{I} \times \mathbf{I}; X)\) 的拓扑由度量 \(\delta\)——由 \(\delta(h, h') = \sup_{u \in \mathbf{I} \times \mathbf{I}} d(h(u), h'(u))\) 给出——定义，且空间对这个度量是完备的（TG, X, p. 20, 推论及 TG, X, p. 9, 推论 1）；由 TG, X, p. 24, 定理 1，这个拓扑是可数型的。于是空间 \(\mathcal{C}_{c}(\mathbf{I} \times \mathbf{I}; X)\) 是波兰空间。空间 Z 亦然，因为 K 是波兰空间（TG, IX, p. 57, 命题 1）。

设 H 是由满足下列条件的三元组 \((\varepsilon,\varepsilon^{\prime},h)\) 组成的 Z 的子集：h 是连接 \(c_{\varepsilon}\) 与 \(c_{\varepsilon^{\prime}}\) 的一个严格同伦。从 Z 到 \(X^{2}\) 中的映射——由 \(a_{t}\colon(\varepsilon,\varepsilon^{\prime},h)\mapsto(c_{\varepsilon}(t),h(t,0))\)、\(b_{t}\colon(\varepsilon,\varepsilon^{\prime},h)\mapsto(c_{\varepsilon^{\prime}}(t),h(t,1))\) 与 \(c_{t}\colon(\varepsilon,\varepsilon^{\prime},h)\mapsto(h(0,t),h(1,t))\) 给出——对每个 \(t\in I\) 都连续，因为映射 \(\varepsilon\mapsto c_{\varepsilon}\)（从 K 到 \(\Omega_{a}(X)\)）连续，映射 \(h\mapsto h(s,t)\)（从 \(\mathcal{C}_{c}(\mathbf{I}\times\mathbf{I};\mathbf{X})\) 到 X 中）亦然。按定义，H 是诸集合 \(a_{t}^{-1}(\Delta_{\mathrm{X}})\)、\(b_{t}^{-1}(\Delta_{\mathrm{X}})\) 与 \(c_{t}^{-1}((a,a))\)（对 \(t\in I\)）之交，其中 \(\Delta_{X}\) 表示 X 的对角线。这就证明了 H 是 Z 的一个闭子集。

于是，H 是波兰空间。设 \(p \colon \mathrm{Z} \to \mathrm{K} \times \mathrm{K}\) 是典范投影；按定义我们有 \(\Gamma = p(\mathrm{H})\)。由于 \(\mathrm{K} \times \mathrm{K}\) 是分离的，\(\Gamma\) 是 \(\mathrm{K} \times \mathrm{K}\) 的一个苏斯林子空间（TG, IX, p. 59, 定义 2）。由 IV, p. 362, 命题 7，它是 \(\mathrm{K} \times \mathrm{K}\) 的一个可达子集。

假设 \(\Gamma\) 不贫。空间 K × K 是紧拓扑空间，因而是贝尔空间（TG, IX, p. 55, 定理 1）。给群 \(G_{0}\)（集合 \(\{0,1\}\) 的置换所成的群）赋以离散拓扑，并让乘积拓扑群 \(G = G_{0}^{\mathbf{N}}\) 对角地作用在 \(K = \{0,1\}^{\mathbf{N}}\) 上。于是群 G 通过映射 (g, (x,y)) ↦ (x, g · y) 连续地作用在 K × K 上。由引理 4，由元素 \(g \in G\)（满足 \(\Gamma \cap g \cdot \Gamma \neq \varnothing\)）所成的集合 V 是 G 的单位元的一个邻域。

设 \(g \in V\)；设 \((\varepsilon, \varepsilon') \in \Gamma \cap g \cdot \Gamma\)。由于 \((\varepsilon, \varepsilon') \in \Gamma\)，有 \(\varepsilon\,R\,\varepsilon'\)；由于 \((\varepsilon, \varepsilon') \in g \cdot \Gamma\)，有 \(g^{-1} \cdot (\varepsilon, \varepsilon') = (\varepsilon, g^{-1}\cdot\varepsilon') \in \Gamma\)，因而 \(\varepsilon\,R\,g^{-1}\cdot\varepsilon'\)。我们这样就证明了：对每个 \(g \in V\)，存在 \(\varepsilon \in K\)，使 \(\varepsilon\) 与 \(g\varepsilon\) 对 R 等价。

对 \(m \in \mathbf{N}\)，记 \(\tau_m\) 为 G 的这样的元素：它的所有项都等于 \(e\)，只有指标为 \(m\) 的那一项等于非平凡元素 \(\tau\)（属于 \(G_0\)）。存在一个整数 \(m\) 使 \(\tau_m\) 属于 V；于是选取 \(\varepsilon \in K\)，使 \(\varepsilon\) 与 \(\varepsilon' = \tau_m \cdot \varepsilon\) 对 R 等价。这蕴含圈 \(c_{\varepsilon}\) 与 \(c_{\varepsilon'}\) 严格同伦。按构造，这两个圈在区间 \([0, 2^{-m-1}]\) 与 \([2^{-m}, 1]\) 上一致；在区间 \([2^{-m-1}, 2^{-m}]\) 上，一个是像为 a 的常值映射，另一个是映射 \(t \mapsto c_{m}(2^{m+1}t - 1)\)。由此得出 \(c_{m}\) 严格同伦于像为 \(\{a\}\) 的常值圈，从而得到矛盾。引理 5 由此得证。

我们现在来完成定理 2 的证明。由命题 8，存在一个连续单射映射 \(\gamma\)（从 \(\{0,1\}^{\mathbf{N}}\) 到 K 中），其像与 R 的每个等价类至多交于一点。若 \(k\) 与 \(k'\) 是 \(\{0,1\}^{\mathbf{N}}\) 的两个不同元素，则圈 \(c_{\gamma(k)}\) 与 \(c_{\gamma(k')}\) 在 X 中不严格同伦，且映射 \(\{0,1\}^{\mathbf{N}} \to \pi_1(\mathrm{X},a)\)（由 \(k \mapsto [c_{\gamma(k)}]\) 给出）是单射的。特别地，\(\operatorname{Card}(\pi_1(\mathrm{X},a)) \geqslant \operatorname{Card}(\{0,1\}^{\mathbf{N}}) = \operatorname{Card}(\mathfrak{P}(\mathbf{N}))\)。由于 X 是可数型的可度量化拓扑空间，\(\Omega_a(\mathrm{X})\) 亦然（TG, X, p. 24, 定理 1）。于是，\(\Omega_a(\mathrm{X})\) 同胚于 \([0,1]^{\mathbf{N}}\) 的一个子空间（TG, IX, p. 18, 命题 12），且

\[
\begin{array}{r l} & {\mathrm{Card} (\Omega_ {a} (\mathbf {X})) \leqslant \mathrm{Card} ([ 0, 1 ] ^ {\mathbf {N}}) = \mathrm{Card} (\mathfrak {P} (\mathbf {N}) ^ {\mathbf {N}})} \\ & {\qquad = \mathrm{Card} (\mathfrak {P} (\mathbf {N} \times \mathbf {N})) = \mathrm{Card} (\mathfrak {P} (\mathbf {N})).} \end{array}
\]

更有，\(\operatorname{Card}(\pi_{1}(\mathrm{X},a))\leqslant\operatorname{Card}(\mathfrak{P}(\mathbf{N}))\)。于是由 E, III, p. 25 的推论 2 得出，\(\pi_{1}(\mathrm{X},a)\) 具有连续统的基数，此即所要证明的。

例. --- 设 P 是这样的拓扑空间：它是诸圆周——以 \((2/n,0)\) 为中心、过平面 \(R^{2}\) 的原点（\(n \geqslant 1\)）——之并（III, p. 336, 习题 6）。P 的庞加莱群具有连续统的基数（III, p. 338, 习题 9）。

